\documentclass{../note}

\begin{document}

\begin{zadanie}{440}
Które z następujących stwierdzeń są prawdziwe dla dowolnego porządku $(X, \le)$ i dowolnych $A, B \subseteq X$? A jeśli ten porządek jest liniowy?
\begin{enumerate}[label=(\alph*)]
    \item Jeśli istnieje $\sup(A \cup B)$, to istnieją $\sup A$ i $\sup B$?
    \item Jeśli istnieją $\sup A$ i $\sup B$, to istnieje $\sup(A \cup B)$?
    \item Jeśli istnieją $\sup A$, $\sup B$ i $\sup(A \cup B)$, to $\sup(A \cup B) = \sup\{\sup A, \sup B\}$?
\end{enumerate}
\end{zadanie}
\begin{rozwiazanie}
\begin{enumerate}[label=(\alph*)]
    \item Niech $A = [0, 1), B = (2, 3]$ oraz $X = A \cup B$ i $\leq$ to relacja mniejsze lub równe na liczbach rzeczywistych, o której wiemy, że jest liniową relacją częściowego porządku. 3 jest jedynym ograniczeniem górnym, więc $\sup(A\cup B) = 3$ istnieje. Jeśli $x \in X$ jest ograniczeniem górnym, to $x \in B$, ale wtedy $\frac{x+2}{2} < x$ też jest ograniczeniem górnym, czyli $x$ nie jest supremum $A$. Zatem $\sup A$ nie istnieje.
    \item Bez wymagania na liniowość możemy wziąć taki kontrprzykład: $X = \{a, b\}, A = \{a\}, B = \{b\}$ oraz każdy element jest w relacji tylko z samym sobą. Wtedy $\sup A = a$, $\sup B = b$, ale $\sup(A \cup B) = \sup\{a, b\}$ nie istnieje, bo nie ma ograniczeń górnych. Jeśli relacja jest liniowa, to niech $a = \sup A, b = \sup B$. Skoro relacja jest liniowa, to jedno z $a, b$ jest większe od drugiego; BSO niech to będzie $a$. Wtedy $\forall_{x\in A} a \geq x$ i $\forall_{x\in B} a \geq b \geq x$, czyli $\forall_{x\in A\cup B} a \geq x$. Jeśli jest jakieś inne ograniczenie górne $c$ na $A \cup B$, to jest to ograniczenie górne na $A$, czyli $c \geq a$. Zatem $a = \sup(A\cup B)$.
    \item 
\end{enumerate}
\end{rozwiazanie}

\begin{zadanie}{622}
Jakiej mocy jest zbiór wszystkich łańcuchów
\begin{itemize}[label=--]
    \item w zbiorze $\mathbb{N} - \{0\}$, uporządkowanym przez relację podzielności?
    \item w zbiorze słów nad alfabetem $\{a, b\}$, uporządkowanym prefiksowo?
    \item w zbiorze słów nad alfabetem $\{a, b\}$, uporządkowanym leksykograficznie?
\end{itemize}
\end{zadanie}
\begin{rozwiazanie}
\begin{itemize}[label=--]
    \item Łańcuchy są zbiorami liczb naturalnych, a zbiorów liczb naturalnych jest $|\mathsf{P}(\N)| = \mathfrak{c}$, więc jest co najwyżej $\mathfrak{c}$ łańcuchów.
\end{itemize}
\end{rozwiazanie}

\end{document}